\section{乘积的坐标环}\label{sec:coordinate-rings-of-products}

在第~\ref{sec:products-of-affine-varieties}~节已证明仿射簇的乘积是仿射簇. 现在可以用两个因子的坐标环来计算这个乘积的坐标环.

\begin{Theorem}{乘积的坐标环}\leavevmode\label{thm:coordinate-ring-of-product-is-tensor-product}
对任意仿射簇 $X$ 和 $Y$, 有典范的 $K$-代数同构
\[
  K[X\times Y]=K[X]\otimes_{K}K[Y].
\]
\end{Theorem}

在给出一般证明之前, 先回到一个具体的例子, 再次明确目标.

\begin{Example}{$\mathbb{A}^{3}$ 中一条抛物线的坐标环}\leavevmode\label{ex:coordinate-ring-of-parabola-in-three-space}
令 $X=\mathcal{V}(y-x^{2})\subseteq\mathbb{A}^{2}$ 且 $Y=\mathcal{V}(z-1)\subseteq\mathbb{A}^{1}$, 于是
\[
  X\times Y=\mathcal{V}(y-x^{2},z-1)\subseteq\mathbb{A}^{3},
\]
这是 $\mathbb{A}^{3}$ 中 $z=1$ 平面内的一条抛物线. 可以 (直接或通过零点定理) 证明
\[
  K[X\times Y]=\frac{K[x,y,z]}{\langle y-x^{2},z-1\rangle},
\]
用关系 $[y]=[x^{2}]$ 和 $[z]=1$ 消去所有出现的 $y$ 和 $z$ 后, 得到
\[
  K[X\times Y]\cong K[x].
\]
另一方面, 分别考虑 $X$ 和 $Y$ 的坐标环, 有
\[
  K[X]=\frac{K[x,y]}{\langle y-x^{2}\rangle}\cong K[x]
  \quad\text{且}\quad
  K[Y]=\frac{K[z]}{\langle z-1\rangle}\cong K.
\]
因此由 Exercise~\ref{ex:tensor-isomorphism-of-multivariate-polynomial-rings} 或 Exercise~\ref{ex:ring-is-tensor-unit} 的结论, 有
\[
  K[X]\otimes_{K}K[Y]\cong K[x]\otimes_{K}K=K[x],
\]
所以 $K[X\times Y]$ 和 $K[X]\otimes_{K}K[Y]$ 确实同构.
\end{Example}

\begin{Proof}\textbf{of Theorem~\ref{thm:coordinate-ring-of-product-is-tensor-product}.}\quad
设 $X$ 和 $Y$ 为仿射簇. 为定义 Theorem~\ref{thm:coordinate-ring-of-product-is-tensor-product} 中的典范同构, 考虑函数
\[
  \begin{gathered}
    \phi:K[X]\times K[Y]\to K[X\times Y],\\
    \phi(F,G)=F\times G,
  \end{gathered}
\]
其中 $F\times G$ 是 $X\times Y$ 上由下式给出的函数:
\[
  (F\times G)(a,b)=F(a)G(b).
\]
如果 $F,G$ 是多项式函数, 则 $F\times G$ 也是 $X\times Y$ 上的多项式函数. 建议读者在 Exercise~\ref{ex:product-of-polynomial-functions-is-polynomial} 中验证这一点. 此外, $\phi$ 是 $K$-双线性的乘法映射, 因而诱导 $K$-代数同态
\[
  \widehat{\phi}:K[X]\otimes_{K}K[Y]\to K[X\times Y].
\]
注意 $\widehat{\phi}$ 确实是典范的. 特别地, 它不依赖于为 $X$ 和 $Y$ 选取在仿射空间中的表示. 只剩证明 $\widehat{\phi}$ 是双射, 由此就可知它是同构.

(满射性) 令 $H\in K[X\times Y]$. 选取表示 $X\subseteq\mathbb{A}^{m}$ 与 $Y\subseteq\mathbb{A}^{n}$, 可将 $H$ 看成多项式函数 $h\in K[x_{1},\ldots,x_{m},y_{1},\ldots,y_{n}]$ 在 $X\times Y$ 上的限制. 写成
\[
  h(x,y)=\sum_{i=1}^{\ell}f_{i}(x)g_{i}(y),
\]
其中 $f_{1},\ldots,f_{\ell}\in K[x_{1},\ldots,x_{m}]$ 和 $g_{1},\ldots,g_{\ell}\in K[y_{1},\ldots,y_{n}]$ 是某些多项式. 令 $F_{1},\ldots,F_{\ell}\in K[X]$ 为限制 $f_{1},\ldots,f_{\ell}$ 得到的多项式函数, 类似地得到 $G_{1},\ldots,G_{\ell}\in K[Y]$. 则
\[
  \widehat{\phi}\left(\sum_{i=1}^{\ell}F_{i}\otimes G_{i}\right)=H,
\]
把两边在 $X\times Y$ 中任意一点 $(a,b)$ 处取值即可看出.

(单射性) 设 $H\in K[X]\otimes_{K}K[Y]$ 且 $\widehat{\phi}(H)=0$, 要论证 $H=0$. 注意, $H$ 可以用许多不同方式写成简单张量之和:
\[
  H=\sum_{i=1}^{\ell}F_{i}\otimes G_{i}.
\]
固定一个使 $\ell$ 最小的此类表达式. 我们将证明 $\ell=0$, 从而 $H=0$.

反设 $\ell>0$. 这意味着 $F_{\ell}\neq0$, 否则就可以省去 $H$ 的表达式中的最后一项, 与 $\ell$ 的最小性矛盾. 因此可选取 $a\in X$ 使 $F_{\ell}(a)\neq0$. 由于 $\widehat{\phi}(H)=0\in K[X\times Y]$, 对任意 $b\in Y$ 有
\[
  0=\widehat{\phi}(H)(a,b)=\sum_{i=1}^{\ell}F_{i}(a)G_{i}(b)
  \quad\Longrightarrow\quad
  0=\sum_{i=1}^{\ell}F_{i}(a)G_{i}\in K[Y].
\]
因为 $F_{\ell}(a)\neq0$, 可以将 $G_{\ell}\in K[Y]$ 写成
\[
  G_{\ell}=-F_{\ell}(a)^{-1}\bigl(F_{1}(a)G_{1}+\cdots+F_{\ell-1}(a)G_{\ell-1}\bigr).
\]
将其代入 $H$ 的表达式, 并用张量积关系化简, 得到
\[
  H=\sum_{i=1}^{\ell-1}\bigl(F_{i}-F_{\ell}(a)^{-1}F_{i}(a)F_{\ell}\bigr)\otimes G_{i},
\]
与 $\ell$ 的最小性矛盾. 这一矛盾蕴含 $\ell=0$, 证明完毕.
\end{Proof}

最后用一个例子进一步说明 Theorem~\ref{thm:coordinate-ring-of-product-is-tensor-product}.

\begin{Example}{$\mathbb{A}^{3}$ 中一条抛物线上的多项式函数}\leavevmode\label{ex:polynomial-functions-on-parabola-in-three-space}
回到 Example~\ref{ex:coordinate-ring-of-parabola-in-three-space} 中的簇 $X\times Y$. 多项式函数 $H\in K[X\times Y]$ 的一个例子是
\[
  \begin{gathered}
    H:X\times Y\to K,\\
    H(a,b,c)=ac+b^{2},
  \end{gathered}
\]
这是多项式 $h(x,y,z)=xz+y^{2}$ 的限制. 若 $F_{1},F_{2}\in K[X]$ 分别是多项式 $f_{1}(x,y)=x$ 和 $f_{2}(x,y)=y^{2}$ 的限制, 而 $G_{1},G_{2}\in K[Y]$ 分别是多项式 $g_{1}(z)=z$ 和 $g_{2}(z)=1$ 的限制, 则
\[
  \widehat{\phi}(F_{1}\otimes G_{1}+F_{2}\otimes G_{2})=H.
\]
当然, 回忆 $X\times Y=\mathcal{V}(y-x^{2},z-1)$, 也可以将 $H$ 写作
\[
  H(a,b,c)=a+a^{4},
\]
这时有
\[
  \widehat{\phi}\bigl((F_{1}+F_{1}^{4})\otimes1\bigr)=H.
\]
$K[X]\otimes_{K}K[Y]$ 中这两个看似不同的元素都映到 $H$, 反映出它们实际上并不不同:
\[
  F_{1}\otimes G_{1}+F_{2}\otimes G_{2}=(F_{1}+F_{1}^{4})\otimes1
    \in K[X]\otimes_{K}K[Y],
\]
可以由 $X$ 和 $Y$ 的定义方程以及张量积关系验证这一点.
\end{Example}

\BookExerciseHeading
\begin{BookExercises}
  \item\label{ex:product-of-polynomial-functions-is-polynomial} 设 $X,Y$ 为仿射簇, $F\in K[X]$ 和 $G\in K[Y]$ 为多项式函数. 证明函数
  \[
    \begin{aligned}
      F\times G:X\times Y&\to K,\\
      (a,b)&\mapsto F(a)G(b)
    \end{aligned}
  \]
  是 $X\times Y$ 上的多项式函数.
  \item\label{ex:coordinate-ring-of-cylinder} 圆柱面 $\mathcal{V}(x^{2}+y^{2}-1)\subseteq\mathbb{A}^{3}$ 的坐标环是什么? 分别将答案写成张量积和 $K[x,y,z]$ 的商环.
  \item\label{ex:coordinate-ring-of-product-with-finite-set} 设 $X\subseteq\mathbb{A}^{m}$ 是仿射簇, 且 $Y\subseteq\mathbb{A}^{n}$ 是含 $r$ 个点的集合.
  \begin{enumerate}
    \item 从几何上描述 $X\times Y\subseteq\mathbb{A}^{m}\times\mathbb{A}^{n}$.
    \item 证明 $K[Y]\cong K^{r}$, 并利用 Exercise~\ref{ex:tensor-distributes-over-direct-sum} 和~\ref{ex:ring-is-tensor-unit} 推出 $K[X\times Y]$ 同构于 $r$ 份 $K[X]$ 的直和.
    \item 解释 (a) 中的几何陈述与 (b) 中的代数陈述之间的关系.
  \end{enumerate}
  \item\label{ex:pullback-of-product-projection} 设 $X,Y$ 为仿射簇, 且 $P:X\times Y\to X$ 为投影映射 $P(a,b)=a$. 将 $P^{*}$ 显式描述为同态
  \[
    K[X]\to K[X]\otimes_{K}K[Y].
  \]
  \item\label{ex:pullback-of-pair-of-polynomial-maps} 设 $X,Y,Z$ 为仿射簇, $F:Z\to X$ 和 $G:Z\to Y$ 为多项式映射. 定义函数
  \[
    \begin{aligned}
      H:Z&\to X\times Y,\\
      c&\mapsto(F(c),G(c)).
    \end{aligned}
  \]
  \begin{enumerate}
    \item 证明 $H$ 是多项式映射.
    \item 将 $H^{*}$ 显式描述为同态
    \[
      K[X]\otimes_{K}K[Y]\to K[Z].
    \]
  \end{enumerate}
  \item\label{ex:three-axes-coordinate-ring-not-nontrivial-tensor} 证明 $K$-代数 $K[x,y,z]/\langle xy,xz,yz\rangle$ 不能写成张量积 $A\otimes_{K}B$, 其中 $A$ 和 $B$ 是均不等于 $K$ 的 $K$-代数.
  \item\label{ex:coordinate-rings-of-finite-products} 陈述 Theorem~\ref{thm:coordinate-ring-of-product-is-tensor-product} 对有限个仿射簇乘积的推广, 并用归纳法证明一般陈述.
  \item\label{ex:tensor-product-of-domains-depends-on-ground-field} 假设 $K$ 是代数闭的, 且 $A,B$ 为有限生成的既约 $K$-代数.
  \begin{enumerate}
    \item 将 Proposition~\ref{prop:product-of-irreducible-varieties-is-irreducible} 与已学过的其他结果结合, 说明若 $A$ 和 $B$ 都是整环, 则 $A\otimes_{K}B$ 也是整环.
    \item 通过论证作为 $\mathbb{R}$-代数有
    \[
      \mathbb{C}\otimes_{\mathbb{R}}\mathbb{C}\cong\mathbb{C}\oplus\mathbb{C},
    \]
    而 $\mathbb{C}$ 是整环, $\mathbb{C}\oplus\mathbb{C}$ 不是, 证明当 $K$ 不是代数闭的时, (a) 的结论可能不成立.
  \end{enumerate}
  \item\label{ex:dimension-of-product-from-noether-bases} 论证 $K[X]$ 和 $K[Y]$ 的 Noether 基给出
  \[
    K[X\times Y]=K[X]\otimes_{K}K[Y]
  \]
  的一个 Noether 基, 从而给出 Proposition~\ref{prop:dimension-additive-on-products} 的另一种证明.
\end{BookExercises}
